\chapter{Carotid Endarterectomy}
\section*{Overview}
Carotid endarterectomy removes atherosclerotic plaque from a carotid artery to improve cerebral blood flow and reduce stroke risk in selected patients.
\section*{Why This Test or Procedure Is Ordered}
It is considered for significant symptomatic carotid stenosis and selected asymptomatic high-grade stenosis after evaluation of surgical risk and medical therapy.
\section*{How This Test or Procedure Is Performed}
Through a neck incision, the artery is opened, plaque is removed, and the artery is closed directly or with a patch. A temporary shunt may maintain blood flow during the repair.
\section*{Risks and Recovery}
Stroke, heart attack, cranial-nerve injury, bleeding, infection, and recurrent narrowing are possible. Blood-pressure control and long-term vascular risk reduction remain important.
\section*{References}
\begin{enumerate}\item MedlinePlus. Carotid Artery Surgery. U.S. National Library of Medicine; 2025.\item Current Medical Diagnosis and Treatment. McGraw-Hill; 2025.\end{enumerate}
\clearpage
